\documentclass{article}

% PACKAGES

% correct encoding and typography
\usepackage[english]{babel}
\usepackage[utf8]{inputenc}
\usepackage[T1]{fontenc}

% math
\usepackage{amsfonts, latexsym, mathtools, amsthm, amssymb, amscd}
\usepackage{euscript}
\usepackage{esvect}
\usepackage{bm}
\usepackage{physics}

% aesthetics color options
\usepackage{xcolor, graphicx}
\usepackage{pdfpages}

% to access the last page
\usepackage{lastpage}

% better document formatting
\usepackage{fancyhdr}
\usepackage[a4paper, right=1in, left=1in, bottom=1.2in, top=1.2in, centering]{geometry}
\usepackage{fullpage}

\usepackage{datetime2}

% better stylistic formatting
\usepackage{parskip}
\usepackage{enumitem}
\usepackage{framed}
\usepackage{longtable}

% embedding and referencing
\usepackage{hyperref}
\hypersetup{colorlinks=true,citecolor=blue,urlcolor =black,linkbordercolor={1 0 0}}

% TikZ (workflow diagrams)
\usepackage{tikz}
\usetikzlibrary{arrows.meta,positioning,fit,calc,backgrounds,matrix}

%%%%%%%%%%%%%%%%%%%%%%%%%%%%%%%%%%%%%%%%%%%%%%%%%%%%%%%%%%%%%%%%%%%%%%%%%%%%%%%%%%%%%%%%%%%%%%%%%%%%%%%%%%%%%%%%%%
% Command formatting for ease-of-life

% new math symbols taking no arguments
\newcommand{\kk}{\Bbbk}
\newcommand{\minus}{\smallsetminus}
\newcommand{\dotp}{\boldsymbol{\cdot}}

%redefined math symbols taking no arguments
\newcommand{\<}{\langle}
\renewcommand{\>}{\rangle}
\renewcommand{\iff}{\Leftrightarrow}
\renewcommand{\implies}{\Rightarrow}

% Arrows
\newcommand{\from}{\leftarrow}
\newcommand{\ffrom}{\longleftarrow}
\newcommand{\goesto}{\rightsquigarrow}
\newcommand{\into}{\hookrightarrow}

% Basic bbface (blackboard bold) lettering command:
\newcommand{\bbb}[1]{\ensuremath{\mathbb{#1}}}

% Boldface lettering command used for vectors:
\newcommand{\bvec}[1]{\ensuremath{\mathbf{#1}}}

% Other lettering
\newcommand{\eps}{\varepsilon}


% some operators and delimiters
\DeclareMathOperator{\proj}{proj}
\DeclareMathOperator{\im}{im}
\DeclareMathOperator{\coker}{coker}
\DeclareMathOperator{\coim}{coim}
\DeclareMathOperator{\Span}{Span}
\DeclareMathOperator{\Stab}{Stab}
\DeclareMathOperator{\Orb}{Orb}
\DeclareMathOperator{\Cov}{Cov}
\DeclareMathOperator{\Var}{Var}
\DeclareMathOperator{\EE}{E}
\DeclareMathOperator{\Tor}{Tor}
\DeclareMathOperator{\Ext}{Ext}


% Other more involved shortcuts

%for overline math
\newcommand{\ol}[1]{{\overline{#1}}}

%redefined math symbols taking arguments
\renewcommand{\mod}[1]{\ (\mathrm{mod}\ #1)}

%for easy 2 x 2 matrices
\newcommand{\twobytwo}[1]{\left[\begin{array}{@{}cc@{}}#1\end{array}\right]}

%for easy column vectors of size 2
\newcommand{\tworow}[1]{\left[\begin{array}{@{}c@{}}#1\end{array}\right]}

%%%%%%%%%%%%%%%%%%%%%%%%%%%%%%%%%%%%%%%%%%%%%%%%%%%%%%%%%%%%%%%%%%%%%%%%%%%%%%%%%%%%%%%%%%%%%%%%%%%%%%%%%%%%%%%%%%
%Below are the theorem, definition, example, lemma, etc. body types.

\newtheorem{theorem}{Theorem}
\newtheorem*{proposition}{Proposition}
\newtheorem{lemma}[theorem]{Lemma}
\newtheorem{corollary}[theorem]{Corollary}
\newtheorem{conjecture}[theorem]{Conjecture}
\newtheorem{postulate}[theorem]{Postulate}
\theoremstyle{definition}
\newtheorem{defn}[theorem]{Definition}
\newtheorem{example}[theorem]{Example}

\theoremstyle{remark}
\newtheorem*{remark}{Remark}
\newtheorem*{notation}{Notation}
\newtheorem*{note}{Note}

%%%%%%%%%%%%%%%%%%%%%%%%%%%%%%%%%%%%%%%%%%%%%%%%%%%%%%%%%%%%%%%%%%%%%%%%%%%%%%%%%%%%%%%%%%%%%%%%%%%%%%%%%%%%%%%%%%
%Some formatting tidbits for margins, paragraphs, and removing orphans

% make widows and orphans rare
\clubpenalty =10000
\widowpenalty =10000

% formatting of paragraphs and separation
\setlength{\parindent}{0pt}
\setlength{\parskip}{5pt plus 1pt}
\setlength{\headheight}{13.6pt}

\setlength\marginparwidth{2.2in}
\setlength\marginparsep{1mm}

%%%%%%%%%%%%%%%%%%%%%%%%%%%%%%%%%%%%%%%%%%%%%%%%%%%%%%%%%

\title{Implementation and Ramifications of TAMER-OP: a Toolkit for Algebraic Module Encodings over $\mathbb{R}^n$ and Other Posets}
\author{Erik Mendes Novak}

\begin{document}
	\maketitle
	
	\tableofcontents
	
	
	\section{Introduction}
	
	%%%%%%%%%%%%%%%%%%%%%%%%%%%%%%%%%%%%%%%%%%%%%%%%%%%%%%%%%%%%%%%%%%%%%%%%%%%%%%%%
	% Section 2: Background mathematical information
	%%%%%%%%%%%%%%%%%%%%%%%%%%%%%%%%%%%%%%%%%%%%%%%%%%%%%%%%%%%%%%%%%%%%%%%%%%%%%%%%
	
	\section{Background mathematical information}\label{sec:background}
	
	Persistent homology assigns to data a family of homology groups indexed by a
	(possibly multi-dimensional) parameter, together with linear maps relating the
	groups as the parameter increases.  The resulting algebraic object is a
	\emph{persistence module}.  In one parameter, persistence modules admit a clean
	structure theory (barcodes) and robust metrics; in multiple parameters, the same
	classification program becomes substantially more subtle.  This section records
	the minimal language needed to state the objects studied in this thesis and to
	motivate the computational strategy pursued by TAMER-OP.
	
	Throughout, fix a coefficient field $\kk$ unless explicitly stated otherwise.
	When we discuss implementation choices later, we will emphasize why exact
	arithmetic over $\bbb{Q}$ is often preferred for certifying ranks and homological
	invariants (cf.\ \S\ref{sec:background:conventions}).
	
	\subsection{Posets as categories; persistence modules as functors}
	\label{sec:background:posets-as-cats}
	
	\subsubsection*{Posets and thin categories}
	
	\begin{defn}[Poset]\label{def:poset}
		A \emph{partially ordered set} (poset) is a set $Q$ equipped with a binary relation
		$\leq$ that is reflexive ($q\leq q$), antisymmetric ($q\leq q'$ and $q'\leq q$
		implies $q=q'$), and transitive ($q\leq q'$ and $q'\leq q''$ implies $q\leq q''$).
	\end{defn}
	
	\begin{defn}[Poset as a category]\label{def:poset-as-category}
		Given a poset $(Q,\leq)$, define a (small) category, which we also denote by $Q$,
		whose objects are elements of $Q$, and such that for $q,q'\in Q$ there is a unique
		morphism $q\to q'$ if and only if $q\leq q'$, and no morphism otherwise.  Composition
		is forced by transitivity, and identities are forced by reflexivity.
	\end{defn}
	
	A category arising from a poset in this way is \emph{thin}: between any two objects
	there is at most one morphism.  This viewpoint is the standard bridge between order
	theory and functorial algebraic constructions.
	
	\subsubsection*{Modules over a poset}
	
	\begin{defn}[Poset module / persistence module]\label{def:Q-module}
		Let $(Q,\leq)$ be a poset and let $\mathbf{Vect}_{\kk}$ denote the category of
		$\kk$-vector spaces and $\kk$-linear maps.  A \emph{$Q$-module} (or \emph{persistence
			module indexed by $Q$}) is a covariant functor
		\[
		M \colon Q \longrightarrow \mathbf{Vect}_{\kk}.
		\]
		Concretely, this is the data of:
		\begin{enumerate}[label=(\roman*)]
			\item a vector space $M(q)$ for each $q\in Q$, and
			\item for each comparable pair $q\leq q'$, a linear \emph{structure map}
			\[
			M(q\leq q') \colon M(q)\to M(q'),
			\]
		\end{enumerate}
		subject to the functoriality conditions
		\[
		M(q\leq q)=\mathrm{id}_{M(q)}\qquad\text{and}\qquad
		M(q\leq q'')=M(q'\leq q'')\circ M(q\leq q')\ \text{for }q\leq q'\leq q''.
		\]
	\end{defn}
	
	The indexing poset $Q$ will typically be $\bbb{R}$ or $\bbb{R}^n$ with the usual
	order(s), or $\bbb{Z}^n$ with the product (coordinatewise) order.  The degree $q$
	is sometimes called a \emph{parameter value}, and $M(q)$ a \emph{fiber} or
	\emph{stalk}.
	
	\begin{defn}[Morphisms]\label{def:morphism-of-Q-modules}
		A morphism of $Q$-modules $f\colon M\to N$ is a natural transformation of functors.
		Equivalently, it is a family of linear maps $f_q\colon M(q)\to N(q)$ indexed by
		$q\in Q$ such that for every $q\leq q'$ the following square commutes:
		\[
		\begin{CD}
			M(q) @>{M(q\leq q')}>> M(q')\\
			@V{f_q}VV @VV{f_{q'}}V\\
			N(q) @>{N(q\leq q')}>> N(q').
		\end{CD}
		\]
	\end{defn}
	
	We write $\mathrm{Hom}_Q(M,N)$ for the $\kk$-vector space of such natural
	transformations (when the fibers are finite-dimensional, $\mathrm{Hom}_Q(M,N)$
	itself is often computable by linear algebra).
	
	\subsubsection*{Functor categories are abelian}
	
	The category of $\kk$-vector spaces is abelian, and functor categories into an
	abelian category remain abelian.  In particular, $Q$-modules admit kernels,
	cokernels, images, and short exact sequences.
	
	\begin{proposition}[Pointwise abelian structure]\label{prop:functor-category-abelian}
		For any small category $\mathcal{C}$, the functor category
		$\mathbf{Vect}_{\kk}^{\mathcal{C}}=\mathrm{Fun}(\mathcal{C},\mathbf{Vect}_{\kk})$
		is abelian.  In particular, for $\mathcal{C}=Q$ a poset category, kernels and
		cokernels are computed \emph{objectwise}: for a morphism $f\colon M\to N$ one has
		\[
		(\ker f)(q)=\ker(f_q),\qquad (\mathrm{coker}\,f)(q)=\mathrm{coker}(f_q),
		\]
		and similarly for images and coimages, with the induced structure maps coming from
		functoriality.
	\end{proposition}
	
	\begin{remark}
		Proposition~\ref{prop:functor-category-abelian} is conceptually important because it
		reduces many categorical constructions on persistence modules to familiar linear
		algebra at each parameter value.  Much of TAMER-OP is built around exploiting this
		pointwise nature after reducing an a priori infinite indexing poset to a finite one.
	\end{remark}
	
	\subsubsection*{Upsets, downsets, and indicator modules}
	
	A recurrent theme (central to Miller's theory of tameness) is that many
	computationally tractable persistence modules can be assembled from very simple
	``indicator'' modules associated to order ideals and filters.
	
	\begin{defn}[Upsets and downsets]\label{def:upset-downset}
		Let $Q$ be a poset.
		\begin{enumerate}[label=(\roman*)]
			\item An \emph{upset} (order filter) is a subset $U\subseteq Q$ such that if
			$q\in U$ and $q\leq q'$, then $q'\in U$.
			\item A \emph{downset} (order ideal) is a subset $D\subseteq Q$ such that if
			$q'\in D$ and $q\leq q'$, then $q\in D$.
		\end{enumerate}
		For $q\in Q$, the \emph{principal upset} and \emph{principal downset} are
		\[
		\uparrow q := \{q'\in Q : q\leq q'\},\qquad
		\downarrow q := \{q'\in Q : q'\leq q\}.
		\]
	\end{defn}
	
	\begin{defn}[Indicator modules]\label{def:indicator-modules}
		Let $U\subseteq Q$ be an upset.  The \emph{upset indicator module} $\kk[U]$ is the
		$Q$-module defined by
		\[
		\kk[U](q) :=
		\begin{cases}
			\kk & q\in U,\\
			0 & q\notin U,
		\end{cases}
		\]
		and for $q\leq q'$ the structure map $\kk[U](q\leq q')$ is the identity
		$\kk\to\kk$ when $q,q'\in U$ and is the unique zero map otherwise.  Likewise, for a
		downset $D\subseteq Q$, define the \emph{downset indicator module} $\kk[D]$ by the
		same formula with $D$ in place of $U$.
	\end{defn}
	
	\begin{remark}
		The closure property in Definition~\ref{def:upset-downset} is exactly what guarantees
		functoriality for the ``identity within the region, zero outside'' convention in
		Definition~\ref{def:indicator-modules}.  For general subsets $S\subseteq Q$, the same
		prescription can fail to define a functor because composition along chains may force
		a map to be zero even when both endpoints lie in $S$.
	\end{remark}
	
	\begin{remark}[Yoneda intuition]\label{rmk:yoneda-principal-upset}
		When $Q$ is viewed as a category, the representable functor
		$\mathrm{Hom}_Q(q,-)$ takes values in sets; its $\kk$-linearization yields precisely
		the principal upset indicator $\kk[\uparrow q]$.  Dually, the corepresentable
		$\mathrm{Hom}_Q(-,q)$ linearizes to $\kk[\downarrow q]$.  For finite $Q$, these
		principal indicator modules behave like indecomposable projectives and injectives
		in the functor category, and they serve as building blocks for resolutions (compare
		with \cite{Mil20}).
	\end{remark}
	
	\subsection{1-parameter persistence}\label{sec:background:1d}
	
	This subsection recalls the classical one-parameter story, both to establish
	notation and to emphasize what breaks in the multiparameter setting.
	
	\subsubsection*{From filtrations to persistence modules}
	
	\begin{defn}[Filtration]\label{def:filtration}
		Let $Q$ be a poset.  A \emph{$Q$-filtration} of a topological space $X$ is a functor
		\[
		\mathcal{F}\colon Q\longrightarrow \mathbf{Top}
		\]
		such that $\mathcal{F}(q)$ is a subspace of $X$ for each $q$, and for $q\leq q'$ the
		morphism $\mathcal{F}(q\leq q')$ is the inclusion $\mathcal{F}(q)\hookrightarrow
		\mathcal{F}(q')$.  Equivalently, $\mathcal{F}$ is an order-preserving assignment of
		subspaces.
	\end{defn}
	
	Applying singular homology with coefficients in $\kk$ yields a persistence module.
	
	\begin{defn}[Persistent homology module]\label{def:persistent-homology-module}
		Given a $Q$-filtration $\mathcal{F}$ and an integer $i\geq 0$, the associated
		$i$th \emph{persistent homology module} is the composite functor
		\[
		H_i(\mathcal{F};\kk)\colon Q \xrightarrow{\ \mathcal{F}\ } \mathbf{Top}
		\xrightarrow{\ H_i(-;\kk)\ } \mathbf{Vect}_{\kk}.
		\]
	\end{defn}
	
	The case $Q=(\bbb{R},\leq)$ (or $Q=\bbb{Z}$) is the classical \emph{one-parameter}
	setting.  The filtration may arise, for example, from sublevel sets of a function
	$f\colon X\to \bbb{R}$, where $\mathcal{F}(t)=f^{-1}((-\infty,t])$.
	
	\subsubsection*{Interval modules and barcodes}
	
	The exceptional feature of one-parameter persistence is that (under mild finiteness
	hypotheses) modules admit a complete decomposition into interval summands.
	
	\begin{defn}[Interval module]\label{def:interval-module}
		Let $I\subseteq \bbb{R}$ be an interval (possibly infinite or half-open, depending on
		the chosen convention).  The \emph{interval module} $\kk_I$ is the $\bbb{R}$-module
		defined by
		\[
		\kk_I(t)=
		\begin{cases}
			\kk & t\in I,\\
			0 & t\notin I,
		\end{cases}
		\]
		with structure maps $\kk_I(s\leq t)$ equal to the identity on $\kk$ when $s,t\in I$
		and the zero map otherwise.
	\end{defn}
	
	To state the classification theorem, one usually imposes a finiteness condition,
	such as \emph{pointwise finite-dimensional} (p.f.d.), meaning $\dim_{\kk} M(t)<\infty$
	for all $t$, together with an appropriate tameness condition to rule out pathologies
	in the indexing by $\bbb{R}$ (see \cite{CB15,CDSGO16}).
	
	\begin{theorem}[Barcode decomposition, informal statement]\label{thm:barcode}
		Let $M$ be a p.f.d.\ persistence module indexed by $\bbb{R}$ satisfying a standard
		tameness condition (e.g.\ q-tameness).  Then $M$ decomposes as a direct sum of
		interval modules,
		\[
		M \cong \bigoplus_{\alpha\in A} \kk_{I_\alpha},
		\]
		and the multiset of intervals $\{I_\alpha\}$ is uniquely determined by $M$ up to
		reordering.  This multiset is the \emph{barcode} (or persistence diagram) of $M$.
	\end{theorem}
	
	The earliest formulations of this result for discrete indexing and finitely generated
	modules go back to persistent homology in computational topology; see, for example,
	\cite{ZC05,EH10}.  For $\bbb{R}$-indexed modules,
	a definitive representation-theoretic proof is given by Crawley-Boevey
	\cite{CB15}.
	
	\subsubsection*{Stability and distances}
	
	A major reason barcodes are useful in applications is \emph{stability}: small changes
	in input data yield small changes in the output barcode with respect to a natural
	metric.  There are several equivalent (or closely related) metrics in one parameter,
	notably the bottleneck distance on persistence diagrams and the interleaving distance
	on persistence modules; see \cite{CDSGO16,Les15} for
	systematic treatments.
	
	For the purposes of this thesis, we will use one-parameter stability primarily as a
	benchmark: it is an example of a situation where both (i) a complete discrete
	invariant exists (Theorem~\ref{thm:barcode}) and (ii) that invariant interacts well
	with perturbations.  Both features become substantially more complicated in multiple
	parameters.
	
	\subsection{Multiparameter persistence and why it is hard}\label{sec:background:multi-hard}
	
	\subsubsection*{Definition and basic examples}
	
	\begin{defn}[Multiparameter persistence module]\label{def:multi-parameter-module}
		Fix $n\geq 2$.  A (real-parameter) \emph{$n$-parameter persistence module} is a
		$\bbb{R}^n$-module, i.e.\ a functor $M\colon (\bbb{R}^n,\leq)\to\mathbf{Vect}_{\kk}$,
		where the order is coordinatewise:
		\[
		a=(a_1,\dots,a_n)\leq b=(b_1,\dots,b_n)\quad\Longleftrightarrow\quad
		a_i\leq b_i\ \text{for all }i.
		\]
		Likewise, a \emph{lattice} $n$-parameter persistence module is a $\bbb{Z}^n$-module.
	\end{defn}
	
	Such modules arise from \emph{multifiltrations} $\mathcal{F}\colon \bbb{R}^n\to
	\mathbf{Top}$, for instance from a function $f\colon X\to \bbb{R}^n$ via
	$\mathcal{F}(a)=f^{-1}((-\infty,a_1]\times\cdots\times(-\infty,a_n])$, and then
	applying homology as in Definition~\ref{def:persistent-homology-module}.
	
	\subsubsection*{Why barcodes do not generalize directly}
	
	When $n\geq 2$, the classification landscape changes qualitatively.  There is no
	analog of Theorem~\ref{thm:barcode} in which modules decompose uniquely into
	interval-like building blocks (rectangles, axis-aligned boxes, etc.).  One reason is
	representation-theoretic: the category of multiparameter persistence modules (even
	with finite-dimensionality restrictions) contains families of indecomposables
	depending on continuous parameters and exhibits behavior described as ``wild'' in
	the sense of quiver representation theory \cite{CZ09,Oud15}.
	
	At an informal level, ``wild'' means that the isomorphism classification problem is
	at least as complicated as classifying representations of arbitrary finite-dimensional
	algebras.  Consequently, one does not expect a complete discrete invariant such as a
	barcode to exist in general.
	
	\begin{remark}[What replaces barcodes?]\label{rmk:what-replaces-barcodes}
		In multiparameter persistence, one typically trades completeness for computability and
		stability: instead of a single complete invariant, one studies families of computable
		invariants (rank-based, homological, and/or slice-based) that are informative for
		specific tasks and data regimes.  Much of the design philosophy of TAMER-OP is to build
		a flexible pipeline that supports such families in a uniform algebraic framework.
	\end{remark}
	
	\subsubsection*{Invariants that still make sense}
	
	Even though complete classification fails, many useful invariants remain available.
	
	\begin{defn}[Rank invariant]\label{def:rank-invariant}
		Let $M$ be a $Q$-module.  The \emph{rank invariant} of $M$ is the function
		\[
		\rho_M \colon \{(q,q')\in Q\times Q : q\leq q'\}\longrightarrow \bbb{Z}_{\geq 0},
		\qquad
		\rho_M(q,q') := \mathrm{rank}\big(M(q\leq q')\big).
		\]
	\end{defn}
	
	\begin{defn}[Hilbert function / dimension function]\label{def:hilbert-function}
		The \emph{dimension function} (often called the \emph{Hilbert function} in graded
		settings) of a $Q$-module $M$ is
		\[
		h_M \colon Q\to \bbb{Z}_{\geq 0},\qquad h_M(q):=\dim_{\kk} M(q).
		\]
	\end{defn}
	
	\begin{remark}
		For $Q=\bbb{N}^n$ and suitably finite modules, $h_M$ is literally the Hilbert function
		of a multigraded module over a polynomial ring.  For general posets, it is still a
		basic and often stable summary of size.
	\end{remark}
	
	Another robust idea is to reduce to one parameter by restricting along monotone
	maps from a totally ordered set into $Q$.
	
	\begin{defn}[Restriction to a chain / slice]\label{def:slice}
		Let $C$ be a totally ordered set (often $C=\bbb{R}$ or $\bbb{Z}$) and let
		$\gamma\colon C\to Q$ be an order-preserving map.  The \emph{restriction} (or
		\emph{slice}) of a $Q$-module $M$ along $\gamma$ is the $C$-module $\gamma^*M$
		defined by composition:
		\[
		\gamma^*M := M\circ \gamma \colon C\to \mathbf{Vect}_{\kk}.
		\]
		When $C$ is totally ordered and $\gamma$ is injective, the image of $\gamma$ is a
		chain in $Q$.
	\end{defn}
	
	The restrictions $\gamma^*M$ are one-parameter persistence modules, hence admit
	barcodes under the usual hypotheses.  Slicing therefore packages multiparameter data
	into a family of one-parameter summaries; this idea is central in 2-parameter
	visualization tools such as RIVET \cite{RIVET}.
	
	Finally, there are homological invariants arising from resolutions and derived
	functors (Ext, Tor) that encode multi-parameter interactions.  These will be pursued
	systematically later in the thesis, but the guiding idea is already visible in
	Miller's program: introduce a finiteness condition strong enough to guarantee finite
	presentations and resolutions (and thus computable derived invariants), while weak
	enough to encompass common persistence modules arising from data \cite{Mil20}.
	
	\subsection{Two complementary viewpoints}\label{sec:background:viewpoints}
	
	Multiparameter persistence can be approached from at least two complementary
	perspectives.  The computational architecture of TAMER-OP is explicitly designed to
	move between them.
	
	\subsubsection*{Topological viewpoint: multifiltrations}
	
	At the source is topology: data are used to build a multifiltration
	$\mathcal{F}\colon \bbb{R}^n\to\mathbf{Top}$ (or $\bbb{Z}^n\to\mathbf{Top}$), and
	persistent homology applies a functor $H_i(-;\kk)$ to obtain
	\[
	M_i := H_i(\mathcal{F};\kk)\colon \bbb{R}^n\to \mathbf{Vect}_{\kk}.
	\]
	The topological viewpoint is advantageous for:
	\begin{itemize}
		\item motivating what invariants should mean geometrically (birth, death, and
		interactions across parameters);
		\item importing stability statements from the theory of interleavings and perturbation
		bounds; and
		\item making contact with geometric constructions of filtrations (sublevel sets, distance
		functions, Vietoris--Rips and \v{C}ech complexes, etc.; see \cite{EH10}).
	\end{itemize}
	
	\subsubsection*{Algebraic viewpoint: multigraded commutative algebra}
	
	When the indexing poset is discrete and additive, persistence modules are closely
	related to multigraded modules over polynomial rings.
	
	\begin{proposition}[Equivalence with multigraded modules (discrete setting)]
		\label{prop:graded-module-equivalence}
		Let $n\geq 1$ and consider the poset $(\bbb{N}^n,\leq)$ with coordinatewise order.
		Then the category of $\bbb{N}^n$-modules $M\colon \bbb{N}^n\to \mathbf{Vect}_{\kk}$
		is equivalent to the category of $\bbb{N}^n$-graded modules over the polynomial ring
		$\kk[x_1,\dots,x_n]$ with degree-preserving homomorphisms.
	\end{proposition}
	
	\begin{remark}[Sketch of correspondence]
		Given a persistence module $M$, form the graded $\kk$-vector space
		$\bigoplus_{a\in \bbb{N}^n} M(a)$ and let $x_i$ act by the structure maps
		$M(a\leq a+e_i)\colon M(a)\to M(a+e_i)$.  Functoriality ensures that the $x_i$ commute.
		Conversely, a graded $\kk[x_1,\dots,x_n]$-module determines the structure maps by
		multiplication with monomials.  See \cite{CZ09,Oud15} for
		detailed treatments in the persistence context.
	\end{remark}
	
	This algebraic viewpoint brings powerful tools---resolutions, Betti numbers, Hilbert
	series---but also reveals why classification becomes hard when
	$n\geq 2$: $\kk[x_1,\dots,x_n]$ is no longer a PID, and multigraded modules can have
	intricate syzygies and moduli.
	
	\subsubsection*{Finite posets and incidence algebras}
	
	TAMER-OP often works with finite posets (encoding posets) obtained by compressing a
	large or continuous parameter space.  For a finite poset $P$, $P$-modules can be
	viewed as finite-dimensional representations of the incidence algebra of $P$ (or of a
	quiver with commutativity relations determined by $P$).  This provides both:
	\begin{itemize}
		\item a representation-theoretic interpretation of projectives/injectives and
		resolutions (cf.\ Remark~\ref{rmk:yoneda-principal-upset}), and
		\item an algorithmic opportunity: for finite $P$, most constructions reduce to finite
		linear algebra, and many ``global'' constraints can be enforced by working only
		along cover relations (Hasse edges).
	\end{itemize}
	
	\subsection{Survey of partial solutions}
	\label{sec:background:partial-solutions}
	
	Because multiparameter persistence resists complete classification, a large body of
	work focuses on \emph{partial} invariants and practical pipelines.  The following
	four themes recur throughout the literature and motivate the design choices of this
	thesis.
	
	\subsubsection*{(i) Slicing and fibered barcodes}
	
	The restriction construction in Definition~\ref{def:slice} reduces a multiparameter
	module to a one-parameter module, hence to a barcode.  Varying the slice (for
	instance, varying a line of positive slope in $\bbb{R}^2$) yields a family of barcodes
	that can be used for visualization or for defining distances by integrating (or
	taking suprema of) one-parameter distances.  In the 2-parameter case, the
	\emph{fibered barcode} viewpoint underlies the RIVET framework \cite{RIVET}.
	
	From a computational standpoint, slicing has two key advantages:
	\begin{enumerate}[label=(\roman*)]
		\item it leverages mature one-parameter algorithms, and
		\item it provides stable summaries for applications even when the full module is too
		large to manipulate directly.
	\end{enumerate}
	Its main limitation is incompleteness: different modules can have identical slices.
	
	\subsubsection*{(ii) Rank-based invariants and M\"obius inversion}
	
	The rank invariant (Definition~\ref{def:rank-invariant}) is one of the most basic
	multiparameter invariants.  In one parameter, the rank function determines the
	barcode, and the barcode can be recovered from ranks by an inclusion--exclusion
	(M\"obius inversion) formula.  In higher parameters, an analogous inversion yields
	\emph{signed} measures on axis-aligned rectangles (sometimes called \emph{signed
		barcodes} or \emph{rectangle measures}).  These objects can be used as:
	\begin{itemize}
		\item sparse, computable summaries of the rank invariant on a chosen grid, and
		\item feature representations for statistics/learning pipelines (e.g.\ kernels on
		signed measures).
	\end{itemize}
	From the algebraic side, such M\"obius inversions are naturally phrased on products of
	chains, and they fit well with finite encodings (cf.\ (iv) below).
	
	\subsubsection*{(iii) Commutative algebra invariants via resolutions}
	
	Under the equivalence in Proposition~\ref{prop:graded-module-equivalence}, a
	$\bbb{N}^n$-module becomes a multigraded $\kk[x_1,\dots,x_n]$-module, and one can
	study:
	\begin{itemize}
		\item minimal free resolutions and multigraded Betti numbers,
		\item Tor/Ext groups and their algebra structures,
		\item associated primes and socles, and
		\item Hilbert series and local cohomology.
	\end{itemize}
	These invariants detect multi-parameter interactions that rank functions and slices
	can miss.  However, they require a finiteness condition strong enough to ensure that
	resolutions are finite in a computational sense.  For real-parameter modules, the
	usual noetherian hypotheses fail, motivating Miller's theory of \emph{tameness} as a
	replacement \cite{Mil20}.
	
	\subsubsection*{(iv) Geometric/combinatorial encoding strategies}
	
	A guiding principle in this thesis is that many modules of interest are
	\emph{piecewise constant} over the parameter space in a controlled way.  An
	\emph{encoding} compresses the indexing poset $Q$ (often $\bbb{R}^n$ or $\bbb{Z}^n$)
	to a finite poset $P$ so that the module factors through $P$ and all subsequent
	computations may be carried out on the finite object.
	
	In Miller's framework, tameness can be characterized in several equivalent ways,
	including (among others) the existence of a finite encoding and the existence of
	certain finite ``fringe'' presentations by upset and downset indicator modules
	\cite{Mil20}.  The computational moral is that once such a finite encoding
	is constructed, one can safely deploy the abelian and homological algebra of finite
	poset modules (kernels, images, Ext/Tor, resolutions) using finite linear algebra.
	
	\subsection{Brief software landscape}
	\label{sec:background:software-landscape}
	
	This thesis is not primarily a software survey, but it is useful to clarify how
	different toolchains emphasize different parts of the multiparameter landscape.
	We will return to more detailed comparisons in later sections, after the TAMER-OP data
	model and pipeline are established.
	
	\begin{itemize}
		\item \textbf{RIVET.} In the 2-parameter case, RIVET is centered on the fibered
		barcode: it supports interactive exploration of how one-parameter barcodes vary over
		a family of lines, and it provides infrastructure for rank-based visualization and
		related distances \cite{RIVET}.
		
		\item \textbf{multipers.} The multipers ecosystem (as commonly used in applied TDA)
		targets scalable computation of multiparameter summaries and feature maps (often
		slice-based, rank-based, or image/landscape-type constructions).  Its emphasis is
		typically on practical pipelines rather than complete algebraic structure.  (A precise
		feature-by-feature comparison depends on the chosen backends and will be addressed
		later.)
		
		\item \textbf{TAMER-OP.} TAMER-OP is designed around Miller's tame-module viewpoint:
		build or ingest a presentation over a large indexing poset, compress it via a finite
		encoding poset, and then compute a broad family of invariants (rank-based,
		homological/derived, geometric, and statistical) by systematically exploiting abelian
		and homological algebra over the finite encoding.
	\end{itemize}
	
	\begin{remark}[Setting expectations for comparisons]\label{rmk:comparison-expectations}
		Comparisons are most meaningful when the underlying mathematical object is held fixed.
		For example, comparing fibered barcodes is natural when both toolchains work with
		2-parameter modules and slice-based summaries.  Comparing resolution-based invariants
		is meaningful when both toolchains expose (or can emulate) multigraded homological
		algebra.  TAMER-OP's main distinguishing feature is that it treats \emph{finite encodings}
		as first-class objects, so that many invariants can be computed uniformly once the
		encoding is in place.
	\end{remark}
	
	\subsection{Notation and conventions}\label{sec:background:conventions}
	
	\subsubsection*{Coefficients and exactness}
	
	Unless stated otherwise, modules take values in $\mathbf{Vect}_{\kk}$ for a fixed
	field $\kk$.  In implementations, there is a persistent tradeoff between speed and
	certifiability:
	\begin{itemize}
		\item Over floating-point fields (e.g.\ $\bbb{R}$), numerical linear algebra is fast
		but can introduce rank errors near degeneracies.
		\item Over exact fields (notably $\bbb{Q}$), ranks, kernels, and Smith/row-reduction
		type computations are certifiable but may be more expensive.
	\end{itemize}
	TAMER-OP is designed so that exact computations over $\bbb{Q}$ can be used when needed,
	especially for invariants derived from ranks and homological algebra (kernels,
	cokernels, Ext/Tor dimensions).
	
	\subsubsection*{Orders on $\bbb{Z}^n$ and $\bbb{R}^n$}
	
	We use the coordinatewise order on $\bbb{Z}^n$ and $\bbb{R}^n$:
	\[
	a\leq b \iff a_i\leq b_i\ \text{for all }i.
	\]
	When writing $a<b$, we mean $a\leq b$ and $a\neq b$.  The standard basis vectors are
	$e_1,\dots,e_n$.
	
	For a finite poset $P$, we generally denote its order by $\leq$ as well; the intended
	meaning should be clear from context.  When $P$ is used as an encoding poset for an
	ambient parameter space, elements of $P$ are sometimes called \emph{regions}.
	
	\subsubsection*{Birth/death terminology and (co)presentations}
	
	When working with indicator modules (Definition~\ref{def:indicator-modules}) and their
	linear combinations, we follow the ``birth/death'' intuition from persistence:
	\begin{itemize}
		\item Upsets model \emph{birth} regions (where generators become present and persist
		upward in the order).
		\item Downsets model \emph{death} constraints (where relations become enforced downward
		in the order).
	\end{itemize}
	A \emph{presentation} will typically mean a cokernel of a map between free objects
	built from upset indicators, while a \emph{copresentation} dually involves downset
	indicators.  In Miller's terminology, \emph{fringe presentations} mix these two
	directions and are tailored to tame modules over general posets \cite{Mil20}.
	
	% ============================================================
	% Section 3: Tame modules
	% ============================================================
	
	\section{Tame modules}
	\label{sec:tame-modules}
	
	Multiparameter persistence is often presented as a story about functoriality: a
	multifiltration of spaces (or chain complexes) gives rise to a functor from a parameter
	poset into vector spaces, and this functor is the algebraic object one studies.
	The difficulty is that, once the parameter poset has dimension $\ge 2$, there is no
	reasonable discrete classification of all such functors, even under strong finiteness
	conditions such as pointwise finite-dimensionality; in representation-theoretic language,
	the classification problem becomes ``wild'' in general \cite{CZ09}.
	For computational purposes one therefore needs a \emph{hypothesis of controlled
		complexity}---a notion of finiteness that is strong enough to guarantee that modules admit
	finite descriptions and are stable under the constructions one cares about, but flexible
	enough to encompass the examples arising in applications.
	
	In this thesis we adopt Ezra Miller's framework of \emph{tame modules over posets}
	\cite{Mil20}.  The guiding idea is that the module should have only finitely many
	distinct ``regions of algebraic behavior'' in parameter space, so that it can be
	represented by \emph{finite data} after passing through a suitable \emph{encoding}.
	A primary goal of TAMER-OP is to turn this philosophy into concrete data structures and
	algorithms; the present section sets up the mathematical backbone that makes such an
	implementation sensible.
	
	\subsection{Why ``tame'' is the right hypothesis for computation}
	\label{subsec:why-tame}
	
	The classical $1$-parameter theory enjoys a powerful structure theorem:
	under mild hypotheses, persistence modules decompose as direct sums of interval modules,
	and the resulting multiset of intervals is encoded by a barcode or persistence diagram.
	This viewpoint underlies much of applied persistent homology, as well as practical
	algorithms for computing invariants from filtered data \cite{ZC05}.
	One rigorous formulation of this classification is due to Crawley--Boevey:
	pointwise finite-dimensional persistence modules over a totally ordered index set admit
	interval decompositions under appropriate finiteness hypotheses \cite{CB15}.
	
	The multiparameter setting behaves fundamentally differently.  Even for
	$\bbb{Z}^n$-indexed modules with $n\ge 2$, indecomposables occur in families and the
	category does not admit a discrete complete invariant analogous to a barcode \cite{CZ09}.
	From a computational perspective, this means that one should not expect a \emph{complete}
	classification of multiparameter persistence modules by a finite combinatorial object.
	Instead, one aims for a framework in which
	\begin{enumerate}[label=(\roman*)]
		\item the objects of interest admit \emph{finite representations},
		\item the relevant algebraic constructions (kernels/cokernels, derived functors, etc.)
		stay within the class, and
		\item the invariants one computes are \emph{functorial} and stable under natural
		operations on the input data.
	\end{enumerate}
	
	Miller's notion of tameness is designed with exactly these desiderata in mind.
	Very roughly, a module is tame if it is constant on finitely many regions of parameter
	space and therefore can be expressed via a finite ``model poset'' that records the
	comparability relations among these regions \cite{Mil20}.
	The definition is arranged so that the category of tame modules (with an appropriate
	restriction on morphisms) is abelian, hence supports homological algebra; moreover,
	tameness admits several equivalent characterizations---combinatorial, topological, and
	homological---that are captured by the syzygy theorem (Theorem~\ref{thm:syzygy} below).
	
	\subsection{Encoding posets and encoding maps}
	\label{subsec:encoding-posets}
	
	The central finiteness notion used throughout this thesis is \emph{finite poset encoding}.
	Intuitively, an encoding replaces a large (typically infinite) poset $Q$ by a smaller
	poset $P$ whose elements represent ``regions'' of the parameter space, such that the
	module on $Q$ is obtained by pullback from a module on $P$.
	
	\begin{defn}[Poset encoding {\cite[Def.~4.1]{Mil20}}]
		\label{def:poset-encoding}
		Let $Q$ be a poset, viewed as a category in the usual way.
		An \emph{encoding} of a $Q$-module $M$ by a poset $P$ consists of
		\begin{enumerate}[label=(\roman*)]
			\item an order-preserving map (equivalently, a functor) $\pi \colon Q \to P$, and
			\item a $P$-module $H \colon P \to \mathbf{Vect}_{\kk}$
		\end{enumerate}
		together with an isomorphism of $Q$-modules
		\[
		M \ \cong \ \pi^*H \ := \ H\circ \pi .
		\]
		The encoding is called \emph{finite} if $P$ has finitely many elements and each stalk
		$H(p)$ is a finite-dimensional $\kk$-vector space.
	\end{defn}
	
	\begin{remark}[Encodings as ``finite models'']
		\label{rem:encoding-fibers}
		If $M\cong \pi^*H$, then $M$ is constant on each fiber $\pi^{-1}(p)$: for any
		$q,q'\in Q$ with $\pi(q)=\pi(q')=p$, the only morphism $p\to p$ in the poset category
		is the identity, so the structure maps of $M$ restrict to identities on the shared stalk
		$H(p)$.  Thus a finite encoding exhibits $M$ as having finitely many ``constant
		regions'' indexed by $P$.
		Conversely, in many geometric situations a partition of the parameter space into finitely
		many constant regions determines a finite encoding by taking $P$ to be the poset of
		regions ordered by reachability (or by closure relations); see \cite{Mil20} for the
		general equivalence among the various constructions.
	\end{remark}
	
	The functorial viewpoint is especially important in what follows, so we record the basic
	formalism.
	
	\begin{defn}[Pullback along a poset morphism]
		\label{def:pullback}
		Let $\pi\colon Q\to P$ be a poset morphism.  Precomposition defines a functor
		\[
		\pi^* \colon \mathbf{Vect}_{\kk}^{P} \longrightarrow \mathbf{Vect}_{\kk}^{Q},
		\qquad
		H \longmapsto H\circ \pi.
		\]
		We refer to $\pi^*$ as the \emph{pullback} (or restriction) functor along $\pi$.
	\end{defn}
	
	\begin{lemma}[Exactness of pullback]
		\label{lem:pullback-exact}
		For any poset morphism $\pi\colon Q\to P$, the pullback functor $\pi^*$ is exact.
		In particular, it preserves finite limits and colimits (kernels, cokernels, and finite
		direct sums).
	\end{lemma}
	
	\begin{proof}
		The category $\mathbf{Vect}_{\kk}^{P}$ of $P$-modules is a functor category with values
		in $\mathbf{Vect}_{\kk}$, hence limits and colimits are computed pointwise.
		Given a short exact sequence $0\to H'\to H\to H''\to 0$ of $P$-modules, evaluating at
		any $p\in P$ yields a short exact sequence of vector spaces.  Precomposing with $\pi$
		and then evaluating at $q\in Q$ simply evaluates the original sequence at $\pi(q)$, so
		exactness is preserved at every stalk, hence globally.
	\end{proof}
	
	Definition~\ref{def:poset-encoding} formalizes the idea that the module ``factors
	through'' a finite poset.  The following terminology, also standard in \cite{Mil20}, is
	useful when one wishes to emphasize the dependence on a fixed encoding map.
	
	\begin{defn}[Subordinate encoding {\cite[Def.~4.7]{Mil20}}]
		\label{def:subordinate}
		An encoding map $\pi\colon Q\to P$ is said to be \emph{subordinate} to a $Q$-module $M$
		if there exists a $P$-module $H$ such that $M\cong \pi^*H$.
	\end{defn}
	
	\begin{example}[Indicator modules admit trivial finite encodings {\cite[Ex.~4.6]{Mil20}}]
		\label{ex:indicator-encoding}
		Let $U\subseteq Q$ be an upset (upward-closed set).  Then the indicator module $\kk[U]$
		(defined in Subsection~\ref{subsec:fringe-presentations}) admits a finite encoding by the
		two-element chain $0<1$ via the map $\pi\colon Q\to\{0<1\}$ sending $U$ to $1$ and
		$Q\setminus U$ to $0$; the corresponding module $H$ has $H(1)=\kk$ and $H(0)=0$.
		Dually, any downset indicator module $\kk[D]$ is encoded by the same chain with the
		roles of $0$ and $1$ reversed.
	\end{example}
	
	\subsection{Definitions of tameness in this framework}
	\label{subsec:def-tameness}
	
	The previous subsection defined what it means for a module to be encoded by a fixed
	poset $P$.  Tameness is the assertion that such an encoding exists with $P$ finite.
	
	\begin{defn}[$P$-tame and tame modules]
		\label{def:tame-module}
		Fix a poset $Q$.
		\begin{enumerate}[label=(\roman*)]
			\item Let $P$ be a finite poset.  A $Q$-module $M$ is \emph{$P$-tame} if there exists a
			poset morphism $\pi\colon Q\to P$ and a $P$-module $H$ with finite-dimensional stalks
			such that $M\cong \pi^*H$.
			\item A $Q$-module $M$ is \emph{tame} if it is $P$-tame for some finite poset $P$.
			Equivalently, $M$ admits a finite poset encoding subordinate to it in the sense of
			Definition~\ref{def:poset-encoding}.
		\end{enumerate}
	\end{defn}
	
	\begin{remark}[Comparison with other finiteness notions]
		\label{rem:other-tameness}
		There are several finiteness conditions in the persistence literature, especially over
		$\bbb{R}^n$ and $\bbb{Z}^n$, including variants of ``constructibility'' and ``$q$-tameness''
		based on finiteness of ranks of structure maps.
		Those conditions are well adapted to stability questions and to certain metric structures,
		but they do not by themselves supply a finite combinatorial model that supports the full
		suite of homological computations one might wish to perform.
		Finite encodings are a \emph{stronger, more combinatorial} hypothesis: they guarantee
		that all information about $M$ is carried by a finite poset $P$ together with finite
		linear algebra on the stalks of a $P$-module.  In particular, tameness is expressly
		designed to enable algorithmic computation in principle \cite{Mil20}.
	\end{remark}
	
	A subtle point is that \emph{not every} $Q$-module homomorphism between tame modules
	should be regarded as ``computably tame''.  Allowing all homomorphisms can take one
	outside the tame world (for example, kernels of arbitrary maps need not be finitely
	encoded).  The correct remedy is to restrict to morphisms compatible with a common
	finite description.
	
	\begin{defn}[Tame morphisms and the tame category {\cite[Def.~4.27, Def.~4.29]{Mil20}}]
		\label{def:tame-morphism}
		Let $M,N$ be tame $Q$-modules.
		A homomorphism $\varphi\colon M\to N$ of $Q$-modules is called \emph{tame} if there
		exists a finite ``constant subdivision'' of $Q$ that is subordinate to both $M$ and $N$,
		such that on each constant region the induced linear map is independent of the chosen
		parameter value (see \cite[Def.~4.27]{Mil20} for the precise formulation).
		The \emph{category of tame $Q$-modules} is the category whose objects are tame
		$Q$-modules and whose morphisms are tame homomorphisms.
	\end{defn}
	
	\begin{theorem}[Tame modules form an abelian category {\cite[Prop.~4.31]{Mil20}}]
		\label{thm:tame-abelian}
		For any poset $Q$, the category of tame $Q$-modules (with tame morphisms) is abelian.
		In particular, it has kernels and cokernels, finite products and coproducts, and every
		monomorphism (resp.\ epimorphism) is a kernel (resp.\ cokernel).
	\end{theorem}
	
	\subsection{Why non-abelianity fears are usually irrelevant in computation}
	\label{subsec:tame-nonabelianity-practice}

	The worry voiced by many users of multiparameter persistence is the following.
	The ambient category $\mathbf{Vect}_{\kk}^{Q}$ of $Q$-modules is abelian, so it is natural
	to form kernels, cokernels, homology, and derived functors there.
	However, the full subcollection of \emph{tame} objects inside $\mathbf{Vect}_{\kk}^{Q}$ is
	not automatically stable under these operations if one allows \emph{arbitrary} natural
	transformations between tame objects.  In other words, tameness can fail to be preserved
	if one performs abelian constructions in $\mathbf{Vect}_{\kk}^{Q}$ without restricting the
	class of morphisms.

	This subsection unpacks the mathematical content behind that concern and explains why,
	from the standpoint of computation, the bad behavior is almost never encountered.
	The key point is that there are two distinct notions of ``morphism between tame modules'',
	and only one of them is compatible with finite data structures.

	\subsubsection*{Homomorphisms of tame objects versus tame morphisms}

	Recall from Remark~\ref{rem:tame-morphisms-practice} that a finite encoding
	$\pi\colon Q\to P$ converts a tame $Q$-module into a $P$-module for a finite poset $P$.
	In that finite setting, morphisms are ordinary natural transformations in
	$\mathbf{Vect}_{\kk}^{P}$, hence are controlled by finitely many linear maps.

	Inside $\mathbf{Vect}_{\kk}^{Q}$, however, a natural transformation
	$\varphi\colon M\to N$ between tame modules can vary on infinitely many degrees even when
	$M$ and $N$ are each constant on finitely many regions.  Miller isolates the maps that
	are compatible with finite encodings by requiring \emph{constancy of the transformation
	on constant regions}:
	a homomorphism $\varphi$ is called \emph{tame} if there exists a finite constant subdivision
	subordinate to both $M$ and $N$ such that on each region $I$ the composite map
	$M_I\to M_q\to N_q\to N_I$ is independent of $q\in I$ \cite[Def.~4.27]{Mil20}.
	The category of tame modules uses these tame homomorphisms as its morphisms
	\cite[Def.~4.29 and Rem.~4.30]{Mil20}.  The resulting category is abelian by
	Theorem~\ref{thm:tame-abelian}.

	\subsubsection*{A concrete failure mode}

	The distinction between arbitrary homomorphisms and tame morphisms is not cosmetic.
	The following example demonstrates that kernels can leave the tame world if the morphism
	is allowed to vary infinitely within a constant region.

	\begin{example}[A kernel of a homomorphism of tame modules need not be tame {\cite[Ex.~4.25]{Mil20}}]
		\label{ex:kernel-not-tame}
		Let $Q=\mathbb{R}^2$ with the coordinatewise partial order, and let
		$L=\{(x,y)\in\mathbb{R}^2 : x+y=1\}$ be the antidiagonal line.
		Let $U\subseteq Q$ be the closed upset above $L$, and let $U^{\circ}$ be its interior.
		Consider the tame modules
		\[
			M=\kk[U]\oplus \kk[U],
			\qquad
			N=\kk[U]/\kk[U^{\circ}],
		\]
		where $N$ is supported on $L$.
		Define a surjection $\varphi\colon M\twoheadrightarrow N$ by taking $\varphi$ to be
		zero on $U^{\circ}$ and, for each point $(x,y)\in L$, sending the two standard basis
		vectors of $M_{(x,y)}=\kk^2$ to the scalars $y$ and $-x$ in $N_{(x,y)}=\kk$.

		Then $K=\ker\varphi$ is \emph{not} tame.  Indeed, for $(x,y)\in L$, the stalk $K_{(x,y)}$
		is the one-dimensional subspace of $\kk^2$ spanned by $(x,y)$, and these lines are pairwise
		distinct as $(x,y)$ varies along $L$.  Consequently, any constant subdivision subordinate to $K$
		must separate distinct points of $L$ into distinct constant regions, forcing infinitely many regions.
	\end{example}

	The phenomenon in Example~\ref{ex:kernel-not-tame} is that $M$ and $N$ are each constant on a
	finite subdivision of $Q$ (three regions: below $L$, on $L$, and above $L$), but the map
	$\varphi$ is \emph{not} constant on the region $L$ itself.  Hence $\varphi$ is not a tame
	homomorphism, so there is no reason for its kernel to remain tame.

	\subsubsection*{The tame category is closed under the usual constructions}

	Once morphisms are required to be tame, the basic abelian operations are safe.

	\begin{lemma}[Kernels and cokernels of tame homomorphisms are tame {\cite[Lem.~4.28]{Mil20}}]
		\label{lem:tame-kernel-cokernel}
		Let $\varphi\colon M\to N$ be a tame homomorphism of $Q$-modules.
		Then $\ker\varphi$ and $\coker\varphi$ are tame modules, and the canonical maps
		$\ker\varphi\to M$ and $N\to\coker\varphi$ are tame morphisms.
		The same statement holds in the semialgebraic, PL, subanalytic, or class~$X$ settings
		when the constant subdivision witnessing tameness is of the corresponding type.
	\end{lemma}

	\begin{proof}
		Choose a finite constant subdivision of $Q$ that is subordinate to $M$, $N$, and $\varphi$
		(as required in the definition of a tame homomorphism).
		For each constant region $I$, fix identifications $M_q\cong M_I$ and $N_q\cong N_I$
		for all $q\in I$ so that the structure maps of $M$ and $N$ agree with the transition
		maps among the regionwise vector spaces.
		By tameness of $\varphi$, the induced linear map $M_I\to N_I$ is independent of $q\in I$.

		Define regionwise vector spaces
		\[
			K_I:=\ker(M_I\to N_I),
			\qquad
			C_I:=\coker(M_I\to N_I).
		\]
		Using the structure maps between regions, these assemble into $Q$-modules $K$ and $C$.
		By construction, $K=\ker\varphi$ and $C=\coker\varphi$ in the abelian category
		$\mathbf{Vect}_{\kk}^{Q}$, and the chosen constant subdivision is subordinate to both $K$
		and $C$.  Hence $K$ and $C$ are tame.  The auxiliary hypotheses (PL, semialgebraic, etc.)
		are preserved because the same subdivision witnesses tameness in those settings.
	\end{proof}

	Lemma~\ref{lem:tame-kernel-cokernel} is the formal version of the slogan that
	``computations with tame modules stay tame'' as long as the maps used in the computation
	come from finite data.  Example~\ref{ex:kernel-not-tame} shows that if one insists on allowing
	\emph{all} homomorphisms between tame objects, then kernels and cokernels can escape the tame
	world, but those escaping morphisms are precisely the ones that do not admit a finite encoding.

	\subsubsection*{Why the pathological maps do not appear in algorithms}

	Any explicit computation in this thesis ultimately proceeds by passing to some finite encoding
	$\pi\colon Q\to P$ and working in the finite functor category $\mathbf{Vect}_{\kk}^{P}$.
	In that regime, every morphism is represented by finitely many linear maps, and pulling back
	along $\pi$ produces a tame homomorphism in the sense of Definition~\ref{def:tame-morphism}.
	Exactness of pullback (Lemma~\ref{lem:pullback-exact}) then implies that kernels and cokernels
	computed at the finite level remain finite after pulling back.

	More explicitly, if $M\cong\pi^{*}H$ and $N\cong\pi^{*}G$ for $P$-modules $H$ and $G$, and if
	$\bar\varphi\colon H\to G$ is a morphism in $\mathbf{Vect}_{\kk}^{P}$, then
	$\varphi:=\pi^{*}\bar\varphi$ is a tame homomorphism of $Q$-modules, and
	\[
		\ker(\varphi)\cong \pi^{*}(\ker\bar\varphi),
		\qquad
		\coker(\varphi)\cong \pi^{*}(\coker\bar\varphi),
	\]
	because $\pi^{*}$ is exact.  Therefore the outputs of the standard constructions
	(kernel, cokernel, homology of a complex, and derived functors computed via finite
	resolutions) remain tame whenever the computation is expressed at the level of a finite
	encoding.  In particular, the ``risky'' behavior in Example~\ref{ex:kernel-not-tame} cannot
	occur unless one is willing to specify a morphism that varies on infinitely many degrees
	inside a constant region, which is not compatible with the finite data structures used for
	encoding and computation.

	\subsubsection*{Recent results that restore full abelianity under geometric hypotheses}

	The preceding discussion explains why Miller restricts morphisms in the tame category.
	There is also a complementary line of results showing that, for many geometric classes of
	finite encodings, the restriction is in practice unnecessary: under suitable hypotheses on
	encoding fibers, every homomorphism between encodable modules can be realized inside a
	finite encoding.

	One precise formulation is due to Waas \cite{Waa24}, who introduces a connectivity condition
	on encoding fibers ensuring that any finite set of encodable modules admits a \emph{common}
	encoding that captures \emph{all} morphisms between them \cite[Prop.~3.5]{Waa24}.
	In particular, the resulting class of encodable persistence modules forms a full abelian
	subcategory of $\mathbf{Vect}_{\kk}^{Q}$ \cite[Thm.~1.1]{Waa24}.
	For parameter posets $Q=\mathbb{R}^n$, his hypotheses are satisfied by natural families of
	``geometric'' regions built from topologically closed PL, semialgebraic, or subanalytic
	upsets (and their complements and finite unions).

	Conceptually, the point is that the failure in Example~\ref{ex:kernel-not-tame} relies on
	large disconnected families of degrees (here the continuum of points on $L$) along which a map
	can vary independently.  Connectivity hypotheses on encoding fibers rule out that degree of
	freedom, forcing morphisms to be controlled by finite combinatorial data.  Waas also notes that
	one can alternatively prove the corresponding subanalytic statement by translating the problem
	into constructible sheaf theory and invoking results of Miller and Kashiwara-Schapira
	\cite[Rem.~3.19]{Waa24}, \cite{Mil23,KS18}.

	These developments give a precise mathematical interpretation of the informal claim that
	``if the input is tame, then the output is almost always tame'': the non-tame outputs arise
	from morphisms that do not admit finite encodings, and such morphisms are excluded either
	(i) by working in the tame category (tame morphisms by definition), or
	(ii) by restricting attention to the geometric subclasses of tame modules for which all
	relevant morphisms are automatically encodable.

	\subsection{Fringe presentations as a computable interface}
	\label{subsec:fringe-presentations}
	
	Finite encodings are the most conceptually direct way to impose combinatorial finiteness,
	but they are not always the most convenient format for \emph{input} or for certain
	computations.  Miller's framework supplies a complementary algebraic interface:
	\emph{fringe presentations}.  These generalize the idea of presenting a module via
	generators and relations, where generators are organized as upsets (birth regions) and
	relations as downsets (death regions).
	
	We begin by recalling the relevant combinatorics of subsets of a poset.
	
	\begin{defn}[Upsets and downsets]
		\label{def:upset-downset}
		Let $Q$ be a poset.
		\begin{enumerate}[label=(\roman*)]
			\item An \emph{upset} (or upper set) in $Q$ is a subset $U\subseteq Q$ such that
			$q\in U$ and $q\le q'$ imply $q'\in U$.
			\item A \emph{downset} (or lower set) in $Q$ is a subset $D\subseteq Q$ such that
			$q\in D$ and $q'\le q$ imply $q'\in D$.
		\end{enumerate}
		For $q\in Q$, the \emph{principal upset} generated by $q$ is $\uparrow q := \{q'\in Q : q\le q'\}$,
		and the \emph{principal downset} is $\downarrow q := \{q'\in Q : q'\le q\}$.
	\end{defn}
	
	\begin{defn}[Indicator modules for upsets and downsets]
		\label{def:indicator-modules}
		Let $Q$ be a poset and let $S\subseteq Q$.
		The \emph{indicator module} (or \emph{indicator $Q$-module}) $\kk[S]$ is the
		$Q$-module defined on objects by
		\[
		\kk[S](q) :=
		\begin{cases}
			\kk & \text{if } q\in S,\\
			0   & \text{if } q\notin S,
		\end{cases}
		\]
		and on morphisms by the rule that for $q\le q'$ the structure map
		$\kk[S](q\le q')\colon \kk[S](q)\to \kk[S](q')$ is the identity $\kk\to\kk$ when both
		$q,q'\in S$, and is $0$ otherwise.
		When $S=U$ is an upset, $\kk[U]$ is called an \emph{upset module}; when $S=D$ is a
		downset, $\kk[D]$ is called a \emph{downset module}.
	\end{defn}
	
	\begin{remark}[Projectives and injectives on a finite poset]
		\label{rem:proj-inj-finite-poset}
		When $P$ is a finite poset, the category $\mathbf{Vect}_{\kk}^{P}$ can be identified
		with the category of finite-dimensional representations of the Hasse quiver of $P$ with
		commutativity relations (equivalently, with modules over the incidence algebra of $P$).
		In this setting, principal upset modules $\kk[\uparrow p]$ are the indecomposable
		projective objects and principal downset modules $\kk[\downarrow p]$ are the
		indecomposable injective objects.  This is one reason indicator modules are the correct
		building blocks for computable homological algebra; see \cite{Mil20} for a detailed
		development in the general poset setting.
	\end{remark}
	
	The next definition packages a tame module using finitely many upsets and downsets
	together with a matrix over $\kk$.  This is the abstraction of ``birth'' and ``death''
	data familiar from barcodes, but adapted to general posets.
	
	\begin{defn}[Fringe presentation {\cite[Def.~3.16]{Mil20}}]
		\label{def:fringe-presentation}
		Let $Q$ be a poset and let $M$ be a $Q$-module.
		A \emph{fringe presentation} of $M$ consists of
		\begin{enumerate}[label=(\roman*)]
			\item a direct sum $F = \bigoplus_{i\in I}\kk[U_i]$ of upset modules,
			\item a direct sum $E = \bigoplus_{j\in J}\kk[D_j]$ of downset modules, and
			\item a $Q$-module homomorphism $\varphi\colon F\to E$
		\end{enumerate}
		such that $\im(\varphi)\cong M$ and the component maps
		$\kk[U_i]\to \kk[D_j]$ satisfy a mild connectivity condition (see \cite[Def.~3.14]{Mil20}).
		The fringe presentation is \emph{finite} if the index sets $I$ and $J$ are finite.
	\end{defn}
	
	\begin{remark}[Relation to classical presentations]
		\label{rem:fringe-as-presentation}
		If $Q=\bbb{Z}^n$ with the usual partial order, then $Q$-modules can be interpreted as
		$\bbb{Z}^n$-graded modules over the polynomial ring $\kk[x_1,\dots,x_n]$ (Section~2.4).
		In that language, upset modules behave like ``free/flat'' pieces and downset modules
		behave like ``injective'' pieces; the map $\varphi$ organizes generators (birth upsets)
		and relations (death downsets) in a way that is robust under changes of geometry and
		does not depend on choosing a coordinate-aligned region decomposition.
	\end{remark}
	
	Fringe presentations become especially concrete once one writes the map $\varphi$ in
	coordinates.
	
	\begin{defn}[Monomial matrix {\cite[Def.~3.17]{Mil20}}]
		\label{def:monomial-matrix}
		Let $\varphi\colon F\to E$ be a \emph{finite} fringe presentation with
		$F=\bigoplus_{i=1}^{m}\kk[U_i]$ and $E=\bigoplus_{j=1}^{r}\kk[D_j]$.
		Choosing the natural bases on the $1$-dimensional stalks of each indicator summand,
		the map $\varphi$ can be encoded by an $r\times m$ matrix $\Phi=(\Phi_{j i})$ with
		entries in $\kk$, whose columns are indexed by the birth upsets $U_i$ and whose rows
		are indexed by the death downsets $D_j$.  We refer to $\Phi$ as a \emph{monomial matrix}
		for $\varphi$.
	\end{defn}
	
	\begin{lemma}[Support condition for monomial matrices {\cite[Prop.~3.18]{Mil20}}]
		\label{lem:monomial-support}
		With notation as in Definition~\ref{def:monomial-matrix}, one necessarily has
		$\Phi_{j i}=0$ whenever $U_i\cap D_j=\varnothing$.  Conversely, any matrix
		$\Phi\in \kk^{r\times m}$ satisfying $\Phi_{j i}=0$ when $U_i\cap D_j=\varnothing$
		determines a fringe presentation map $\varphi\colon F\to E$.
	\end{lemma}
	
	\begin{example}[A single interval in one parameter]
		\label{ex:interval-fringe}
		Let $Q=\bbb{R}$ with its total order and consider the interval module $\kk[a,b)$.
		This module can be realized as the image of a map from the upset module $\kk[a,\infty)$
		to the downset module $\kk(-\infty,b)$, and the corresponding monomial matrix is the
		$1\times 1$ matrix $[1]$.  For a direct sum of intervals (a barcode), the matching of
		births to deaths is reflected by a monomial matrix that is a permutation matrix after
		choosing orderings of births and deaths; compare \cite[Ex.~3.19]{Mil20} and the
		interval decomposition theorem \cite{CB15}.
	\end{example}
	
	\subsection{What tameness buys you: mathematical consequences}
	\label{subsec:tame-math}
	
	The principal reason tameness is a useful organizing principle is that it is not a single
	isolated definition but rather part of a web of equivalent finiteness conditions, each
	suited to different tasks (geometric modeling, algebraic presentation, homological
	computation).  The main structural result is Miller's syzygy theorem.
	
	\begin{theorem}[Syzygy theorem {\cite[Thm.~6.12]{Mil20}}]
		\label{thm:syzygy}
		Let $Q$ be a poset and let $M$ be a $Q$-module.  Then $M$ is tame if and only if it
		admits (and hence admits all of) the following finite descriptions:
		\begin{enumerate}[label=(\roman*)]
			\item a finite poset encoding subordinate to $M$;
			\item a finite fringe presentation of $M$;
			\item finite resolutions by upset modules and by downset modules.
		\end{enumerate}
		Moreover, the constructions can be chosen compatibly with suitable auxiliary geometric
		structures (e.g.\ semialgebraic or piecewise-linear settings) when $Q$ sits inside a
		partially ordered real vector space, and tame morphisms lift to morphisms of the
		corresponding finite presentations/resolutions.
	\end{theorem}
	
	\begin{remark}[Why equivalences matter for software]
		\label{rem:equivalences-matter}
		Theorem~\ref{thm:syzygy} has two practical interpretations:
		\begin{enumerate}[label=(\roman*)]
			\item It legitimizes switching among data representations depending on what is being
			computed.  One might build a module from geometric data (suggesting a region-based
			encoding), transform it to a fringe presentation to expose ``birth/death'' structure,
			and then compute derived invariants via an indicator resolution.
			\item It explains why it is sensible to base a computational pipeline on finite encodings.
			If a module is tame, then \emph{some} finite encoding exists, and any other finite
			representation can (in principle) be converted to an encoding suitable for computation.
		\end{enumerate}
		TAMER-OP is designed around these equivalences: the library implements multiple entry points
		(fringe-style and region-style) and reduces computations to finite poset models.
	\end{remark}
	
	A more elementary but extremely important consequence is that, once an encoding is fixed,
	many constructions can be carried out on the finite poset $P$ and then transported back
	to $Q$.
	
	\begin{lemma}[Reduction along an encoding]
		\label{lem:reduction-along-encoding}
		Suppose $M\cong \pi^*H$ where $\pi\colon Q\to P$ is a finite encoding and $H$ is a
		$P$-module.
		Then every structure map of $M$ is obtained by applying $H$ to the image under $\pi$:
		for $q\le q'$ in $Q$,
		\[
		M(q\le q') \ = \ H\!\big(\pi(q)\le \pi(q')\big).
		\]
		In particular, any invariant of $M$ that is defined functorially from the diagram of
		vector spaces and structure maps (for example, ranks of structure maps, limits/colimits,
		or derived functors computed from resolutions) can be computed from $H$ on the finite
		poset $P$.
	\end{lemma}
	
	\begin{proof}
		The identity $M=H\circ\pi$ is the definition of pullback.  The final statement follows
		because $\pi^*$ is exact (Lemma~\ref{lem:pullback-exact}) and because the data of $M$
		is completely determined by the data of $H$ together with $\pi$.
	\end{proof}
	
	\subsection{What tameness buys you: algorithmic consequences}
	\label{subsec:tame-algo}
	
	The abstract consequences above translate directly into a computational strategy.
	
	\paragraph{Finite linear algebra.}
	If $P$ is finite and $H$ is a $P$-module with finite-dimensional stalks, then $H$ is
	specified by:
	\begin{enumerate}[label=(\roman*)]
		\item a finite list of vector spaces $H(p)$, and
		\item linear maps $H(p\le p')$ subject only to functoriality constraints.
	\end{enumerate}
	In a finite poset, functoriality can be enforced by recording maps only along cover
	relations (the Hasse diagram) and composing along paths.  This is precisely the kind of
	finite linear-algebraic data structure on which one can implement kernels, images,
	cokernels, ranks, and derived constructions.
	
	\paragraph{Homological algebra in a finite abelian category.}
	The abelian category $\mathbf{Vect}_{\kk}^{P}$ has enough projectives and injectives,
	and its basic building blocks are the principal upset and downset modules
	(Remark~\ref{rem:proj-inj-finite-poset}).  Consequently, one can compute projective and
	injective resolutions by explicit finite algorithms, and then compute $\Tor$ and $\Ext$
	groups by linear algebra on chain complexes.
	When a $Q$-module is tame, Theorem~\ref{thm:syzygy} guarantees that such finite
	resolutions exist \emph{upstairs} in $Q$ as well (in the appropriate tame category),
	but the key computational reduction is that it suffices to resolve the encoded object on
	a finite poset model.
	
	\subsection{Bridge to implementation}
	\label{subsec:tame-bridge}
	
	The remainder of this thesis is concerned with turning the preceding theory into a
	usable computational toolkit.  The next sections explain:
	\begin{enumerate}[label=(\roman*)]
		\item which concrete finite presentations and encodings TAMER-OP supports,
		\item how finite encoding posets are constructed from geometric or algebraic input,
		\item how invariants and derived functors are computed on finite poset models, and
		\item how these computations feed into a statistics pipeline for data analysis.
	\end{enumerate}
	In short: tameness provides the \emph{mathematical contract} that makes the entire
	data-ingestion-to-statistics pipeline finite and computable.
	
	
	
	\section{Overview of TAMER-OP}
	

	
	
	\subsection{Design goals stated mathematically}
	
	
	
	\subsection{The end-to-end pipeline}
	
	\subsection{Core mathematical objects in the API}
	
	\subsection{Architectural layering}
	
	\subsection{Reproducibility and correctness}
	
	\section{The presentation formats}
	
	\subsection{Philosophy: multiple equivalent front-ends}
	
	\subsection{Finite fringe presentations}
	
	\subsection{$\mathbb{Z}^n$ presentations}
	
	\subsection{$\mathbb{R}^n$ and piecewise-linear presentations}
	
	\subsection{Conversions between formats}
	
	\subsection{Validation and other sanity checks}
	
	\subsection{Serialization and I/O}
	
	\section{Finite encoding constructors}
	
	\subsection{The encoding problem}
	
	\subsection{Region/signature idea}
	
	\subsection{Encoding algorithms by ambient domain}
	
	\subsection{Constructing the finite poset structure}
	
	\subsection{Point location and evaluation map}
	
	\subsection{Building a common encoding for multiple modules}
	
	\subsection{Complexity and failure modes}
	
	\subsection{Correctness and diagnostics}
	
	\section{Invariants}
	
	\subsection{Meta principle: invariants factor through the encoding}
	
	\subsection{Dimension-type invariants}
	
	\subsection{Rank-type invariants}
	
	\subsection{Mobius inversions and signed measures}
	
	\subsection{Slice-based invariants and distances}
	
	\subsection{Vectorizations and kernels}
	
	\subsection{Homological invariants derived from resolutions}
	
	\subsection{Geometry-aware invariants for when $Q = \mathbb{R}^n$}
	
	\subsection{Stability and robustness discussion}
	
	\subsection{Practical API summary}
	
	\section{Homological Algebra and Derived Functors}
	
	\subsection{The Abelian Category of $P$-modules}
	
	\subsection{Projective and Injective Objects}
	
	\subsection{Resolutions}
	
	\subsection{Ext and Tor: definitions and computational model}
	
	\subsection{Multiplicative structures and higher structure}
	
	\subsection{Double complexes and spectral sequences}
	
	\subsection{The ergonomics of derived categories}
	
	\subsection{Relationship to multigraded commutative algebra}
	
	\subsection{What you use all this for in TDA}
	
	\section{Transport and comparison}
	
	\subsection{Change of posets as functorial transport}
	
	\subsection{Comparing modules on different encodings}
	
	\subsection{Comparison via invariants}
	
	\subsection{Interoperability with multipers and RIVET}
	
	\subsection{Case studies}
	
	\section{Statistics}
	
	\subsection{Data model for experiments}
	
	\subsection{Feature extraction recipes}
	
	\subsection{Basic inferential statistics}
	
	\subsection{Learning-oriented workflows}
	
	\subsection{Diagnostic plots}
	
	\subsection{Reproducibility and reporting}
	
	\section{Speed optimizations}
	
	\subsection{Where time is spent}
	
	\subsection{Tradeoffs in exactness versus speed}
	
	\subsection{Sparsity-aware algorithms}
	
	\subsection{Caching and memoization strategies}
	
	\subsection{Parallelism}
	
	\subsection{Geometry backend: controlling combinatorial explosion}
	
	\section{Benchmarks}
	
	\subsection{Philosophy and reproducibility}
	
	\subsection{Microbenchmarks (unit tests)}
	
	\subsection{End-to-end benchmarks}

	\subsection{Comparisons with multipers and RIVET}
	
	\subsection{Discussion of results}
	
	\section{Conclusion}
	aa
	
	% ------------------------------------------------------------
	% References
	% ------------------------------------------------------------

	
	\begin{thebibliography}{99}
		
		\bibitem[Mil20]{Mil20}
		E.~Miller.
		\newblock \emph{Homological algebra of modules over posets}.
		\newblock \href{https://arxiv.org/abs/2008.00063}{arXiv:2008.00063}, 2020.
		
		

		\bibitem[Waa24]{Waa24}
		L.~Waas.
		\newblock \emph{Notes on abelianity of categories of finitely encodable persistence modules}.
		\newblock \href{https://arxiv.org/abs/2407.08666}{arXiv:2407.08666}, 2024.

		\bibitem[Mil23]{Mil23}
		E.~Miller.
		\newblock \emph{Stratifications of real vector spaces from constructible sheaves with conical microsupport}.
		\newblock \emph{Journal of Applied and Computational Topology}, 7(3):473--489, 2023.
		\newblock \href{https://doi.org/10.1007/s41468-023-00112-1}{doi:10.1007/s41468-023-00112-1}.

		\bibitem[KS18]{KS18}
		M.~Kashiwara and P.~Schapira.
		\newblock \emph{Persistent homology and microlocal sheaf theory}.
		\newblock \emph{Journal of Applied and Computational Topology}, 2(1--2):83--113, 2018.
		\newblock \href{https://arxiv.org/abs/1705.00955}{arXiv:1705.00955}.
\bibitem[CB15]{CB15}
		W.~Crawley-Boevey.
		\newblock \emph{Decomposition of pointwise finite-dimensional persistence
			modules}.
		\newblock \emph{Journal of Algebra and Its Applications}, 14(5):1550066, 2015.
		\newblock \href{https://doi.org/10.1142/S0219498815500668}{doi:10.1142/S0219498815500668}.
		
		\bibitem[ZC05]{ZC05}
		A.~Zomorodian and G.~Carlsson.
		\newblock \emph{Computing persistent homology}.
		\newblock \emph{Discrete \& Computational Geometry}, 33:249--274, 2005.
		\newblock \href{https://doi.org/10.1007/s00454-004-1146-y}{doi:10.1007/s00454-004-1146-y}.
		
		\bibitem[CZ09]{CZ09}
		G.~Carlsson and A.~Zomorodian.
		\newblock \emph{The theory of multidimensional persistence}.
		\newblock \emph{Discrete \& Computational Geometry}, 42:71--93, 2009.
		\newblock \href{https://doi.org/10.1007/s00454-009-9176-0}{doi:10.1007/s00454-009-9176-0}.
		
		\bibitem[CEH07]{CEH07}
		D.~Cohen-Steiner, H.~Edelsbrunner, and J.~Harer.
		\newblock \emph{Stability of persistence diagrams}.
		\newblock \emph{Discrete \& Computational Geometry}, 37(1):103--120, 2007.
		\newblock \href{https://doi.org/10.1007/s00454-006-1276-5}{doi:10.1007/s00454-006-1276-5}.
		
		\bibitem[ELZ02]{ELZ02}
		H.~Edelsbrunner, D.~Letscher, and A.~Zomorodian.
		\newblock \emph{Topological persistence and simplification}.
		\newblock \emph{Discrete \& Computational Geometry}, 28:511--533, 2002.
		\newblock \href{https://doi.org/10.1007/s00454-002-2885-2}{doi:10.1007/s00454-002-2885-2}.
		
		\bibitem[Les15]{Les15}
		M.~Lesnick.
		\newblock \emph{The theory of the interleaving distance on multidimensional
			persistence modules}.
		\newblock \emph{Foundations of Computational Mathematics}, 15(3):613--650, 2015.
		\newblock \href{https://doi.org/10.1007/s10208-015-9255-y}{doi:10.1007/s10208-015-9255-y}.
		
		\bibitem[LW15]{LW15}
		M.~Lesnick and M.~Wright.
		\newblock \emph{Interactive visualization of 2-D persistence modules}.
		\newblock \href{https://arxiv.org/abs/1512.00180}{arXiv:1512.00180}, 2015.
		
		\bibitem[RIVET]{RIVET}
		M.~Lesnick and M.~Wright.
		\newblock \emph{RIVET: The rank invariant visualization and exploration tool}
		(software).
		\newblock \url{https://github.com/rivetTDA/rivet}.
		
		\bibitem[LS24]{LS24}
		D.~Loiseaux and H.~Schreiber.
		\newblock \emph{multipers: Multiparameter persistence for machine learning}.
		\newblock \emph{Journal of Open Source Software}, 9(103):6773, 2024.
		\newblock \href{https://doi.org/10.21105/joss.06773}{doi:10.21105/joss.06773}.
		
		\bibitem[CDSGO16]{CDSGO16}
		F.~Chazal, V.~de~Silva, M.~Glisse, and S.~Y. Oudot.
		\newblock \emph{The Structure and Stability of Persistence Modules}.
		\newblock SpringerBriefs in Mathematics. Springer, 2016.
		\newblock \href{https://doi.org/10.1007/978-3-319-42545-0}{doi:10.1007/978-3-319-42545-0}.
		
		\bibitem[Oud15]{Oud15}
		S.~Y. Oudot.
		\newblock \emph{Persistence Theory: From Quiver Representations to Data Analysis}.
		\newblock Mathematical Surveys and Monographs, Vol.~209. American Mathematical
		Society, 2015.
		
		\bibitem[EH10]{EH10}
		H.~Edelsbrunner and J.~L. Harer.
		\newblock \emph{Computational Topology: An Introduction}.
		\newblock American Mathematical Society, 2010.
		
		\bibitem[Mac98]{Mac98}
		S.~Mac~Lane.
		\newblock \emph{Categories for the Working Mathematician}.
		\newblock Graduate Texts in Mathematics, Vol.~5. Springer, 2nd edition, 1998.
		\newblock \href{https://doi.org/10.1007/978-1-4757-4721-8}{doi:10.1007/978-1-4757-4721-8}.
		
		\bibitem[Wei94]{Wei94}
		C.~A. Weibel.
		\newblock \emph{An Introduction to Homological Algebra}.
		\newblock Cambridge Studies in Advanced Mathematics, Vol.~38. Cambridge
		University Press, 1994.
		
		\bibitem[Eis95]{Eis95}
		D.~Eisenbud.
		\newblock \emph{Commutative Algebra with a View Toward Algebraic Geometry}.
		\newblock Graduate Texts in Mathematics, Vol.~150. Springer, 1995.
		\newblock \href{https://doi.org/10.1007/978-1-4612-5350-1}{doi:10.1007/978-1-4612-5350-1}.
		
		\bibitem[MS05]{MS05}
		E.~Miller and B.~Sturmfels.
		\newblock \emph{Combinatorial Commutative Algebra}.
		\newblock Graduate Texts in Mathematics, Vol.~227. Springer, 2005.
		\newblock \href{https://doi.org/10.1007/b138602}{doi:10.1007/b138602}.
		
	\end{thebibliography}
	
	
\end{document}